% Lösungen Einheit 3 von 3 – Faktorisieren (zusammenbau v0.1)
\einheitenkopf{Einheit 3 von 3 · Faktorisieren}

%% TODO Lösung der Erklärzeile 17g) fehlt in der Bank
\erg{17}{a) ja – Wurzeln $x$ und $8$; $2 \cdot x \cdot 8 = 16x$ \quad b) $(x + 6) \cdot (x - 6)$ \quad c) $(x + 9) \cdot (x - 9)$ \quad d) $(x + 4) \cdot (x - 4)$ \quad e) $(x + 1) \cdot (x - 1)$ \quad f) $(x + 11) \cdot (x - 11)$ \quad g) \ldots \quad h) $(x + 4)^2$, Mittelglied $2 \cdot x \cdot 4 = 8x$ passt \quad i) $(x - 9)^2$ \quad j) $(4x + 3) \cdot (4x - 3)$ \quad k) $2 \cdot (x^2 - 16) = 2 \cdot (x + 4) \cdot (x - 4)$ \quad l) Nein: Summe zweier Quadrate; die dritte Formel braucht ein Minus dazwischen, die erste und zweite ein Mittelglied. \quad m) $(x - 8)^2 = x^2 - 16x + 64$ – stimmt. \quad n) $(x + 12) \cdot (x - 12) = 0$, also $x_1 = -12$, $x_2 = 12$ \quad o) $(x + 4)^2 - 13$, denn $x^2 + 8x + 16 - 16 + 3$ \quad p) $3 \cdot (x^2 - 16) = 3 \cdot (x + 4) \cdot (x - 4)$; Probe: $3 \cdot (x^2 - 16) = 3x^2 - 48$}
\erg{18}{a) Fehler: eine Summe zweier Quadrate ohne Mittelglied ist keine binomische Formel. Richtig: Der Term lässt sich so nicht faktorisieren. \quad b) Die dritte Formel ergibt eine Differenz, die erste und zweite ein Mittelglied; eine Summe zweier Quadrate ohne Mittelglied entsteht bei keiner.}
